% ============================================================
% §7 工程实施建议 (diangong winning_patterns #9)
% 竞赛: 电工杯 B题 (2026) | Stage 08 产物
% ============================================================

\section{工程实施建议}

基于Q1人口预测、Q2选址优化、Q3定价补贴与Q4鲁棒性分析的全链条定量证据，本文面向地方政府与养老服务运营商提出以下三条可操作的分阶段实施建议。

\subsection{建议一：近期（1年内）按109万元"全覆盖最优方案"启动建设}

Q2的确定性MILP求解（randomSeed=42）在120万元预算约束内找到了覆盖率100\%（10/10小区）的最优方案：F（大型，45万元）+ H（中型，32万元）+ J（中型，32万元），总建设成本109万元，剩余预算11万元。三站"全员满负荷"（利用率92--96\%）的运营状态虽然在$S_2$响应满意度上形成了普遍的0.600悬崖效应，但实现了"不落一户"的全覆盖——这是地方政府最核心的政治目标。建议在首年以109万元的确定性方案为工程招标基准，将剩余的11万元预算用于：(1) F站办公楼租赁补贴（大型站需约200m$^2$，年均租金约5--8万元）；(2) 服务人员岗前培训（约3万元/年）。从社会福利视角看，以109万元公共投资撬动10个小区6,864名老人的基础养老服务可及性，人均建设成本仅为159元/人（远低于行业平均的300--500元/人），具有极高的公共财政效率。

\subsection{建议二：中期（1--3年）建立"浮动补贴率+满意度梯度改善"双机制}

Q3分析揭示了两项关键的结构性政策洞见：(1) 现行"规模定额补贴帽"导致大型站的有效补贴率仅0.73元/次（效率损失63.5\%），而中型站为0.79--0.81元/次——形成"服务越多、补贴效率越低"的逆向激励。建议将补贴上限从"按站点规模定额"改为"按实际服务人次浮动补贴率"，基础方案为"1.5元/次$\times$日非紧急服务人次"，同时设置站点级日补贴上限为"理论补贴额的75\%"以保留财政审慎性。按此方案，F站的有效补贴率将从0.73升至1.13元/次（+55\%），全系统年补贴从223.2万元增至约330万元——增加的107万元财政支出对应的是"消除逆向激励、恢复F站服务意愿"的制度性收益。(2) Q2方案中三站$S_2$全员触底（0.600）的满意度悬崖，可通过追加13万元预算将H站从中型升级为大型——日容量从2,000增至3,000人次，H站的利用率从94.5\%降至63.0\%，其$S_2$将从0.600跃升至1.000（满意度+67\%），B/E/H三小区的综合满意度$S$从0.840升至0.960。建议在第二年度预算周期中优先审批该13万元升级支出，将"全员满负荷"的脆弱平衡升级为"一大多中+适度冗余"的稳健配置。

\subsection{建议三：远期（3--5年）构建"季度监测--预警--触发的动态容量管理闭环}

Q4 Scenario C（银发海啸：半失能与失能转化率各+30\%）的极端压力测试表明，若失能人口额外增长12.7\%，全系统月需求将从247,060次增至约270,000次——超过当前三站总月容量（$7,000\times30=210,000$人次）达28.6\%，当前的全覆盖方案将因容量不足而崩溃。而在成本通胀20\%的情景下（Scenario B），系统通过集约化（从3站压缩为2大站，C+J）仍维持了80\%的覆盖率与0.905的满意度——系统的自适应集约化能力通过了压力测试。

建议建立以季度为频率的老年人口状态监测机制，设置两级预警阈值：
\begin{enumerate}[label=\textbf{黄灯预警：}, leftmargin=*]
    \item 若连续两个季度的半失能$\to$失能转化率超过11\%（基线10\%），则触发"银发海啸黄灯预警"，启动以下预案：(a) 将F站的日容量从3,000临时扩展至3,300（+10\%，通过延长服务时间1小时实现）；(b) 启动E小区附近预留地块的简化审批程序。
\end{enumerate}
\begin{enumerate}[label=\textbf{红灯预警：}, leftmargin=*]
    \item 若失能转化率超过12.5\%（基线+25\%），则触发"容量红灯预警"，立即启动H站扩建（中型$\to$大型，追加投资13万元）的加急招标程序，并同步向市级财政申请应急扩容专项拨款。
\end{enumerate}

这一"监测--预警--触发"的动态闭环，将Q4的灵敏度分析从学术练习转化为可操作的工程风险管理工具——使政府能够在"确定性最优方案"与"不确定性的实物期权预留"之间取得动态平衡。
